%! TEX program = xelatex
%! TEX root = lecture2_slide.tex

\title{\texorpdfstring{$\lambda$}{λ}-Calculus}
\subtitle{Simple Types and their Extensions}
\begin{document}

\bgroup
  \usebackgroundtemplate{\includegraphics[width=\paperwidth]{banner.pdf}}
  \begin{frame}[plain,noframenumbering]
    \maketitle
  \end{frame}
\egroup

%\section{Properties of \texorpdfstring{$\lambda$}{λ}-Calculus}
%
%\begin{frame}
%  \begin{example}
%    Suppose $t$ is a term and $y \not\in \FV(t)$. 
%    Then, consider 
%    \[
%      (\lambda y.\, t)\, ((\lambda x.\, x\,x)(\lambda x.\, x\, x))
%    \]
%  \end{example}
%  \mode<presentation>{\vfill}
%  Observations:
%  \begin{itemize}
%    \item Some evaluation may diverge while some may converge.
%    \item Full $\beta$-reduction lacks for determinacy. 
%  \end{itemize}
%  Question:
%  \begin{itemize}
%    \item Does every path give the same evaluation?
%  \end{itemize}
%\end{frame}
%
%\begin{frame}{Confluence}
%\begin{theorem}[Church-Rosser]
%  Given $u_1$ and $u_2$ with $t \reduce u_1$ and $t\reduce u_2$, there is $v$
%  such that $u_1 \reduce v$ and $u_2 \reduce v$, pictorially,
%  \[
%    \xymatrix{
%      & t \ar@{->>}[rd]^{\beta} \ar@{->>}[ld]_{\beta} \\
%      u_1 \ar@{->>}[rd]_{\beta} & & u_2 \ar@{->>}[ld]^{\beta} \\
%      & v
%    }
%  \]
%\end{theorem}
%
%No matter which way we choose we can always find a \alert{confluent}
%term. 
%
%\end{frame}
%
%\begin{frame}{Normal form}
%  \begin{definition}
%    $M$ is \emph{in normal form} if 
%    there is no $N$ such that $M \onereduce N$, abbreviated to $M \notonereduce$.
%    
%  \end{definition}
%  \begin{lemma}
%    Suppose that $M$ is in normal form. Then 
%      $M \reduce N$ implies $M =_\alpha N$.
%  \label{lem:normal-no-reduction}%
%  \end{lemma}
%  \begin{proof}
%    By induction on the derivation $d$ of $M \reduce N$.
%    \begin{enumerate}
%      \item If $d$ is given by (0-step), then
%        $M \reduce N$ where $M =_\alpha N$ by definition.
%
%      \item If $d$ is given by (n+1-step), then
%        $M \onereduce M'$ and $M' \reduce N$ are derivable for some $M'$. By
%        assumption $M \onereduce N$ is not derivable for any $N$, so by
%        contradiction the statement follows.
%    \end{enumerate}
%  \end{proof}
%
%\end{frame}
%
%\begin{frame}{Corollaries of confluence}
%  \begin{corollary}[Uniqueness of normal forms]
%    Let $M$ be a term with $M \reduce N_1$ and $M\reduce N_2$ where
%     $N_i$'s are in normal form. Then, $N_1 =_\alpha N_2$.
%    \label{coro:uniqueness-normal}%
%  \end{corollary}
%  
%
%  \begin{corollary}[Computationally equal terms have a confluent term]
%    If $M =_\beta N$, then there exists $L$ satisfying
%  \[
%    \xymatrix{
%      M \ar@{->>}[rd]_{\beta} \ar@{}[rr]|{=_\beta} & & N \ar@{->>}[ld]^{\beta} \\
%      & L
%    }
%  \]
%  \label{coro:computational-meanting-confluence}
%  \end{corollary}
%  \begin{proof}[Proof sketch]
%    By induction on the derivation of $M =_\beta N$. 
%  \end{proof}
%
%%\begin{corollary}
%%  Let $=_\beta$ denote the congruence closure of $\onereduce$. 
%%  \begin{enumerate}
%%    \item If $M =_\beta N$, then there exists $L$ such that 
%%      \[
%%        M \reduce L
%%        \; _{\beta*}\longleftarrow N
%%      \]
%%    \item If in addition $N$ is in normal form, then 
%%      $M \reduce N$. 
%%  \end{enumerate}
%%\end{corollary}
%%\begin{block}{Homework}
%%  Show this corollary.
%%\end{block}
%\end{frame}


\section{Simply Typed \texorpdfstring{$\lambda$}{λ}-Calculus: Introduction}

\begin{frame}{Adding types to a language}
  While $\lambda$-calculus is expressive and computationally powerful, it is rather painful to write programs inside $\lambda$-calculus.

  Function can be applied to an arbitrary term which can represent a Boolean value, a number, or even a function, so as a programming language it is not easy to see the \alert{intention} of a program.

  Therefore, we will consider a formal definition of a \alert{typing judgement}
  \[
    \Gamma \vdash t : A
  \]
  which specifies the type $A$ of a term $t$ under a list of free (typed) variables, allowing us to \emph{restrict the formation} of a valid term by typing.
\end{frame}

\section{Simply Typed \texorpdfstring{$\lambda$}{λ}-Calculus: Statics}

%\begin{frame}{Typing judgement}
%  A \alert{typing judgement} is of the form
%  \[
%    \Gamma \vdash t : A
%  \]
%  saying the \emph{term $t$ is of type $\sigma$ under the context $\Gamma$}
%  where 
%  \begin{description}
%    \item[$\Gamma$] is a list of variables $x_i : A_i$ available in $t$
%    \item[$A$] is the type of $t$.
%  \end{description}
%
%  For example, the typing judgement
%  \[
%    x_1:A_1, x_2: A_2 \vdash x_1 : A_1
%  \]
%  means `the term $x_1$ is of type $A_1$ under the context $x_1:A_1, x_2:A_2$'.
%\end{frame}



\begin{frame}{Higher-order function type}
Assume $\mathbb{V}$ is a set of type variables different from variables in untyped $\lambda$-terms. 
(And suppress its existence from now on.)

\begin{definition}
  The judgement $A : \type$ is defined inductively as follows.
  \begin{multicols}{2}
    \begin{prooftree}
      \AXC{\vphantom{$X$}}
      \RightLabel{if $X \in \mathbb{V}$}
      \UIC{$X : \type$}
    \end{prooftree}
    \begin{prooftree}
      \AXC{$A : \type$}
      \AXC{$B : \type$}
      \BIC{$A \to B : \type$}
    \end{prooftree}
  \end{multicols}
  where $A \to B$ represents a function type from $A$ to $B$.
\end{definition}
We say that $A$ is a type if $A : \type$ is derivable.

\end{frame}

\begin{frame}
  The function type is \alert{higher-order}, because
  \begin{enumerate}
    \item functions can be arguments of another function;
    \item functions can be the result of a computation.
  \end{enumerate}
  For example, 
  \begin{description}
    \item[$\;(A_1 \to A_2) \to B$] a function type whose argument is of type $A_1 \to A_2$; 
    \item[$\;A_1 \to (A_2 \to B)$] a function whose return type is $A_2 \to B$. 
  \end{description}
  \vfill
  Following the convention of function application, we introduce the convention for the function type:
  \begin{block}{Convention}
  \[
    A_1 \to A_2 \to \dots A_n \quad  \defeq\quad A_1 \to
    (A_2 \to ( \dots \to (A_{n-1} \to A_n)\dots))
  \]
  \end{block}
\end{frame}

\begin{frame}{Context}
  
\begin{definition}
    A \emph{typing context}~$\Gamma$
    is a sequence
    \[
      \Gamma \equiv x_1 : A_1,\; x_2 : A_2,\; \ldots,\; x_n : A_n
    \]
    of \alert{\emph{distinct variables}} $x_i$ of type $A_i$.
\end{definition}

\begin{definition}
  The membership judgement $\Gamma \ni (x : A)$ is defined inductively:
  
  \begin{multicols}{2}
    \begin{prooftree}
      \AXC{$\vphantom{\Gamma}$}
      \UIC{$\Gamma, x : A \ni x : A$}
    \end{prooftree}
    \begin{prooftree}
      \AXC{$\Gamma \ni x : A$}
      \UIC{$\Gamma, y : B \ni x : A$}
    \end{prooftree}
  \end{multicols}
  We say that $x$ of type $A$ occurs in $\Gamma$ if $\Gamma \ni (x : A)$ if derivable.
\end{definition}
\end{frame}

%\begin{frame}{Hypothetical judgement}
%  A \emph{hypothetical judgement} (or, an inference rule) is of the form
%  \begin{prooftree}
%    \AXC{$J_1$}
%    \AXC{$J_2$}
%    \AXC{$\dots$}
%    \AXC{$J_n$}
%    \RightLabel{(rule)}
%    \QuaternaryInfC{$J$}
%  \end{prooftree}
%  consisting of 
%  \begin{description}
%    \item[hypothesis] a set of judgements $J_i$, for $1 \leq i \leq n$
%    \item[name] the name for the reason why $J_i$'s imply $J$
%    \item[conclusion] a single judgement $J$
%  \end{description}
%  Examples include
%  \begin{enumerate}
%    \item The formation of $\lambda$-terms (var), (app), and (abs).
%    \item The formation of $\alpha$-equivalence $=_\alpha$.
%    \item The formation of simple types $\mathbb{T}$.
%  \end{enumerate}
%\end{frame}

\begin{frame}{Typing rule -- Curry-style typing system}
  The implicit typing system for simply typed $\lambda$-calculus is defined by the following typing rules, i.e.\ 
  inference rules with its conclusion a typing judgement:
  
  \begin{prooftree}
    \AXC{}
    \RightLabel{(var)\quad if $\Gamma \ni (x : A)$\quad }
    \UIC{$\Gamma \vdash_i x : A$}
  \end{prooftree}
  \begin{prooftree}
    \AXC{$\Gamma, x : A \vdash_i t : B$}
    \RightLabel{(abs)}
    \UIC{$\Gamma \vdash_i \lambda x.\; t : A \to B$}
  \end{prooftree}
  \begin{prooftree}
    \AXC{$\Gamma \vdash_i t : A \to B$}
    \AXC{$\Gamma \vdash_i u : A$}
    \RightLabel{(app)}
    \BIC{$\Gamma \vdash_i t\;u : B$}
  \end{prooftree}
  We say that $t$ is a \alert{closed term} if $\vdash t : A$ is derivable. 

  N.B.\ Whether a term $t$ has a typing derivation is a \emph{property} of $t$.
\end{frame}

%\begin{frame}{Derivation}
%  Given a set $\mathcal{R}$ of hypothetical judgements, a \emph{derivation} of
%  \begin{prooftree}
%    \AXC{$J_1$}
%    \AXC{$J_2$}
%    \AXC{$\dots$}
%    \AXC{$J_n$}
%    \QuaternaryInfC{$J$}
%  \end{prooftree}
%  is a tree of instances of rules in $\mathcal{R}$ composed with $J_i$'s as the top-level hypothesises.
%
%  A judgement $J$ is \emph{derivable} (with assumptions $J_i$'s) if there is derivation whose root is $J$ (with assumptions $J_i$'s).
%
%  \begin{example}
%    The fact that
%    \[
%      \lambda x\,y.\, y
%    \]
%    is a $\lambda$-term means the judgement $\lambda x\,y.\, y$ is derivable
%    from the formation rules of $\lambda$-terms without any assumption.
%  \end{example}
%\end{frame}

\begin{frame}{Syntax-directedness}
  A typing system is \emph{syntax-directed} if it has \emph{exactly} one typing rule
  for each term construct.
  \vfill
  By being syntax-directed, every typing derivation can be inverted:
  \begin{lemma}[Typing inversion]
    Suppose that $\Gamma \vdash_i t : A$ is derivable.
    Then,
    \begin{description}
      \item[$t \equiv x$] implies $x : A$ occurs in $\Gamma$.
      \item[$t \equiv \lambda x.\, t'$] implies $A = B \to C$ and $\Gamma, x: B \vdash_i u' : C$.
      \item[$t \equiv u\;v$] implies there is some $A$ such that $\Gamma \vdash_i u : A \to B$ and $\Gamma \vdash_i v : B$.
    \end{description}
  \end{lemma}
  This lemma is particularly useful when constructing a typing derivation by hand.
\end{frame}

\begin{frame}{Typing derivation}
%  The judgement $\vdash_i \lam{x}x : A \to A$, for any type $A$, has a derivation
%  \begin{prooftree}
%    \AXC{}
%    \RightLabel{(var)}
%    \UIC{$x : A \vdash_i x : A$}
%    \RightLabel{(abs)}
%    \UIC{$\vdash_i \lambda x.\; x : (A \to A)$}
%  \end{prooftree}

  For any types $A$ and $B$, the judgement $ \vdash_i \lambda x\,y.\, x : A \to B \to A$
  has a derivation
\begin{prooftree}
  \AXC{}
  \RightLabel{(var)}
  \UIC{$x : A, y: B \vdash_i x : A$}
  \RightLabel{(abs)}
  \UIC{$x : A \vdash_i \lambda y.\, x : B \to A$}
  \RightLabel{(abs)}
  \UIC{$\vdash_i \lambda x\,y.\, x : A \to B \to A$}
\end{prooftree}
Therefore, $\lam{x\,y}x$ is a program of type $A \to B \to A$.

\end{frame}

\begin{frame}{Exercise}
  Derive the typing judgement
    \[
      \vdash_i \lambda f\,g\,x.\, f\,x\, (g\,x) : (A \to B \to C) \to
      (A\to B) \to A \to C
    \]
    for every types $A, B$ and $C$.
\end{frame}

%\begin{frame}{Explicit typing: Typed terms}
%\begin{definition}[Typed terms]
%  The formation $M \; \term_{V}^{\to}$ of typed terms is defined by
%    \begin{prooftree}
%      \AXC{$x \in V$}
%      \UIC{$x \quad \term_{V}^{\to}$}
%    \end{prooftree}
%    \begin{prooftree}
%      \AXC{$M \quad \term_{V}^{\to}$}
%      \AXC{$N \quad \term_{V}^{\to}$}
%      \BIC{$M\, N \quad \term_{V}^{\to}$}
%    \end{prooftree}
%    \begin{prooftree}
%      \AXC{$M \quad \term_{V}^{\to}$}
%      \AXC{$x \in V$}
%      \AXC{$\tau \quad \type$}
%      \TIC{$\lambda {\color{red}x : \tau}.\; M \quad \term_{V}^{\to}$}
%    \end{prooftree}
%\end{definition}
%\end{frame}
%
%\begin{frame}{Explicit typing: Typing rules}
%\begin{definition}[Typing Rules]
%  Typing derivations on \emph{typed terms} are defined by 
%  \begin{prooftree}
%    \AXC{$\Gamma \ni (x : \sigma)$}
%    \RightLabel{(var)}
%    \UIC{$\Gamma \vdash_e x : \sigma$}
%  \end{prooftree}
%  \begin{prooftree}
%    \AXC{$\Gamma \vdash_e M : \sigma \to \tau$}
%    \AXC{$\Gamma \vdash_e N : \sigma$}
%    \RightLabel{(app)}
%    \BIC{$\Gamma \vdash_e M\;N : \tau$}
%  \end{prooftree}
%  \begin{prooftree}
%    \AXC{$\Gamma, {\color{red}x : \sigma} \vdash_e M : \tau$}
%    \RightLabel{(abs)}
%    \UIC{$\Gamma \vdash_e \lambda {\color{red}x : \sigma} .\; M : \sigma \to \tau$}
%  \end{prooftree}
%\end{definition}
%
%\end{frame}
%\begin{frame}{Explicit typing: Unicity}
%
%\begin{proposition}\label{prop:unicity}
%  For every typed term $M$, context~$\Gamma$, and types $\sigma_i$, 
%  \[
%    \Gamma \vdash_e M : \sigma_1
%    \quad\text{and}\quad
%    \Gamma \vdash_e M : \sigma_2
%    \implies
%    \sigma_1 = \sigma_2
%  \]
%\end{proposition}
%\begin{proof}[Proof sketch]
%  Use the inversion lemma and the structural induction on $M$.
%
%  E.g., suppose that $M$ is of the form
%  \[
%    L\;M'
%  \]
%  By inversion there are $\tau_i$ such that $\Gamma \vdash_e L: \tau_i \to
%  \sigma_i$ and $\Gamma \vdash_e M': \tau_i$. By induction hypothesis, $\tau_1 \to
%  \sigma_1 = \tau_2 \to \sigma_2$, so $\sigma_1 = \sigma_2$.
%\end{proof}
%\end{frame}
%
%
%
%\begin{frame}{Type erasure}
%An \emph{erasing map} $|-|\colon \term_{V}^{\to} \to \term_{V}$ is defined by
%  \begin{align*}
%    |x| & = x \\
%    |M\; N| & = |M|\;|N| \\
%    |\lambda x:\sigma.\, M| & = \lambda x.\, |M|
%  \end{align*}
%  \begin{example}
%    \begin{enumerate}
%      \item $|\lambda (f: \sigma \to \tau)\,(x: \sigma).\, f\;x| = \lambda f\, x.\, f\;x$
%      \item $|(\lambda (x: \sigma)\,(y: \tau). y)\;z| = (\lambda x\,y.\, y)\; z$
%    \end{enumerate}
%  \end{example}
%
%$|-|$ is an translation from $\term_{V}^{\to}$ to $\term_{V}$.
%Does $|-|$ respect the behaviour of $\term_{\term_V^{\to}}$?
%\end{frame}
%\begin{frame}{From typed terms to untyped and back}
%\begin{proposition}
%  Let $M$ and $N$ be typed $\lambda$-terms in~$\term_{\lambda_V^{\to}}$. Then, 
%  \begin{align*}
%    \Gamma  \vdash_e M : \sigma & \text{ implies } \Gamma \vdash_i |M| :
%    \sigma \\ 
%    M \reduce N & \text{ implies } |M| \reduce |N|
%  \end{align*}
%\end{proposition}
%
%\begin{proposition}
%  Let $M$ and $N$ be $\lambda$-terms in~$\term_V$. Then, 
%  \begin{enumerate}
%  \item If $\Gamma \vdash_i M : \sigma$, then there is $M' : \term_V^{\to}$ with 
%        $|M'| = M
%        \quad\text{and}\quad
%        \Gamma \vdash_e M' : \sigma$
%      \item If $M \reduce N$ and $M = |M'|$ for some $M' : \term_V^{\to}$,
%      then there exists $N'$ with $|N'| = N$ and $M' \reduce N'$.
%    \end{enumerate}
%\end{proposition}
%\end{frame}

\begin{frame}{Type inference and checking}
  Can we answer the following questions algorithmically?
  \begin{description}
    \item[Type inference] Given a context $\Gamma$ and a term $t$, is there a type $?$
      such that the typing judgement $\Gamma \vdash t : ?$ is derivable? 
    \item[Type checking] Given a context $\Gamma$, a type $A$, and a term $t$, is the typing judgement $\Gamma \vdash t : A$ derivable?
  \end{description}
  
  Typability is reducible to type checking problem of
  \[
    x_0: A \vdash \mathtt{fst} \;x_0\;t : A
  \]

  \begin{theorem}
    Type checking is \emph{decidable} in simply typed $\lambda$-calculus.
  \end{theorem}
  
\end{frame}

\section{Programming in Simply Typed \texorpdfstring{$\lambda$}{λ}-Calculus}

\begin{frame}[allowframebreaks]{Church encodings of natural numbers}
The type of natural numbers is of the form
\[
  \nat_A \defeq (A \to A) \to A \to A
\]
for every type $A$.
  
  \begin{description}
    \item[Church numerals]
      \begin{align*}
        & \bc_n \defeq \lambda f\,x.\,
        f^n x \\
        \vdash{} & \bc_n : \nat_A
      \end{align*}
    \item[Successor]
      \begin{align*}
        & \suc \defeq \lambda n \,f\,x\,.\;f\;(n\;f\;x) \\
        \vdash{} & \suc : \nat_A \to \nat_A
      \end{align*}
    \item[Addition]
      \begin{align*}
        & \add \defeq \lambda n\,m\,f\,x.\; (m\;f)\;(n\;f\;x) \\
        \vdash{} & \add : \nat_A \to \nat_A \to \nat_A
      \end{align*}
    \item[Muliplication] 
      \begin{align*}
        & \mul \defeq \lambda n\,m\,f\,x.\, (m\;(n\;f))\;x\\
      \vdash{} & \mul : \nat_A \to \nat_A \to \nat_A
      \end{align*}
    \item[Conditional]
      \begin{align*}
        & \ifz \defeq \lambda n\,x\,y.\, n\;(\lambda z.\, x)\;y\\
        \vdash {} & \ifz : ?
      \end{align*}
  \end{description}
The type of~$\ifz$ may not be as obvious as you may expect.
Try to find one and justify your guess.

\end{frame}
\begin{frame}{Church encodings of boolean values}
We can also define the type of Boolean values 
for each type variable as
\[
  \bool_A \defeq A \to A \to A
\]
\begin{description}
  \item[Boolean values]
      \[
        \true \defeq \lambda x\,y.\,x 
        \quad\text{and}\quad
        \false \defeq \lambda x\,y.\,y 
      \]
  \item[Conditional]
    \begin{align*}
      & \cond \defeq \lambda b\,x\,y.\, b\,x\,y\\
      \vdash {} & \cond : \bool_{A} \to A \to A \to A
    \end{align*}
\end{description}
\end{frame}

\begin{frame}{Exercise}
  \begin{enumerate}
    \item Define conjunction $\mathtt{and}$, disjunction $\mathtt{or}$, and
      negation $\mathtt{not}$ in simply typed lambda calculus.

    \item Prove that $\mathtt{and}$, $\mathtt{or}$, and $\mathtt{not}$ are well-typed.
  \end{enumerate}
  
\end{frame}


\section{Properties of Simply Typed \texorpdfstring{$\lambda$}{λ}-Calculus}


\begin{frame}[c]{Type safety = Preservation + Progress}

\begin{quote}
  ``Well-typed programs cannot `go wrong'.''\\
  \hfill ---(Milner, 1978)
\end{quote}

\begin{description}
  \item[Preservation] If $\Gamma \vdash t : A$ is derivable and $t \onereduce u$, then $\Gamma \vdash u : A$.
  \item[Progress] If $\Gamma \vdash t : A$ is derivable, then either $t$ is in \emph{normal form} or there is $u$ with $t \onereduce u$.
\end{description}
\vfill
By combing the above two properties, we can extend the progress theorem to $\reduce$: if $\Gamma \vdash t : A$ then $t \reduce u$ for some $\Gamma \vdash u : A$ which is either reducible or in normal form.
\end{frame}

\begin{frame}{Converse of Preservation}
  The converse of preservation might not hold.

\begin{lemma}[Typability of subterms]
  Let $t$ be a term with $\Gamma \vdash t : A$ derivable. Then, for every
  subterm $t'$ of~$t$ there exists $\Gamma'$ such that
  \[
    \Gamma' \vdash t' : A'.
  \]
\end{lemma}
%\begin{proof}
%  By induction on $\Gamma \vdash t : A$.
%\end{proof}
  Recall that 
  \begin{enumerate}
    \item $\mathbf{K}_1 = \lambda x\,y.\, x$
    \item $\Omega = (\lambda x.\, x\,x)\,(\lambda x.\, x\,x)$
  \end{enumerate}
  and $\mathbf{K}_1\;(\lam{x}x)\;\Omega \reduce \mathbf{I}$.

  $\Omega$ is not typable, so $\mathbf{K}_1\,\mathbf{I}\,\Omega$ is not typable.
\end{frame}

\begin{frame}{Preservation Theorem}
    \begin{description}
      \item[Weakening] If $\Gamma \vdash t : A$ and $x \not\in \Gamma$, then
        $\Gamma, x : B \vdash t : A$. 
      \item[Substitution] If $\Gamma, x : A \vdash t : B$ and $\Gamma \vdash u : A$ then $\Gamma \vdash t\subst{u}{x} : B$.
    \end{description}

  \begin{corollary}[Variable renaming]
    If $\Gamma, x : A \vdash t : B$ and $y \not\in \mathrm{dom}(\Gamma)$, then $ \Gamma, y : A \vdash t\subst{y}{x} : B$
    where $\mathrm{dom}(\Gamma)$
    denotes the set of variables which occur in $\Gamma$.
  \end{corollary}

  \begin{theorem}
    For any $t$ and $u$ if $\Gamma \vdash t : A$ is derivable and $t \onereduce u$, then $\Gamma \vdash u : A$. 
  \end{theorem}
  \begin{proof}[Proof sketch]
    By induction on both the derivation of $\Gamma \vdash t : A$ and $t \onereduce u$.
  \end{proof}
  N.B. The only non-trivial case is $\Gamma \vdash (\lam{x}t)\; u : B$ which needs the above results.
%  \begin{proof}
%    $y$ is not in $\Gamma$, so 
%      $\Gamma, y : A, x : A \vdash M$
%    by weakening and by definition $\Gamma, y : A \vdash y : A$.
%    Thus, by substitution, we have
%    \[
%      \Gamma, y : A \vdash M\subst{x}{y} : \sigma
%    \]
%  \end{proof}
\end{frame}


\begin{frame}{Proof of Preservation Theorem}
  \begin{proof}
  By induction on both the derivation of $\Gamma \vdash t : A$ and $t \onereduce u$.

  \begin{enumerate}
    \item Suppose $\Gamma \vdash x : A$.
      However, $x \not\onereduce u$ for any $u$.
      Therefore, it is vacuously true that $\Gamma \vdash u : A$.
    \item Suppose $\Gamma \vdash \lambda x.\, t : A \to B$ and $\lambda x.\, t \onereduce u$.
      Then, $u$ must be $\lambda x.\, u'$ for some $u'$; $\Gamma, x : A \vdash t \colon B$ and $t \onereduce u'$ must be derivable.
      By induction hypothesis, $\Gamma, x : A \vdash u'$ is derivable, so is $\Gamma \vdash \lambda x.\, u'\colon A \to B$.
    \item Suppose $\Gamma \vdash t\; u$.
      Then ...
    \item ...
  \end{enumerate}
\end{proof}
\end{frame}

\begin{frame}{Progress: First attempt}
  \begin{theorem}
    If $\Gamma \vdash t : A$ is derivable, then $t$ is in normal form or there is $u$ with $t \onereduce u$.
  \end{theorem}
  To prove the theorem, we would like to use induction on $\Gamma \vdash t : A$ again.

  However, the fact that $t$ is in normal form does not tell us much what $t$ is.
  Can we characterise $t$ syntactically?
\end{frame}

\begin{frame}{Normal form}
  The notion of normal form can be characterised syntactically:
  \begin{definition}
    Define judgements $\texttt{Neutral}\;t$ and $\texttt{Normal}\;u$ mutually by
    \begin{multicols}{2}
      \begin{prooftree}
        \AXC{$\vphantom{\Gamma}$}
        \UIC{$\texttt{Neutral}\;x$}
      \end{prooftree}
      \begin{prooftree}
        \AXC{$\neutral\;t$}
        \AXC{$\normal\;u$}
        \BIC{$\neutral\;t\;u$}
      \end{prooftree}
      \columnbreak
      \begin{prooftree}
        \AXC{$\neutral\;t$}
        \UIC{$\normal\;t$}
      \end{prooftree}
      \begin{prooftree}
        \AXC{$\normal\;t$}
        \UIC{$\normal\;\lambda x.\,t$}
      \end{prooftree}
    \end{multicols}
  \end{definition}
  \textbf{Idea.} $\neutral\;u$ (resp.\, $\normal\;t$) is derivable
  iff 
  \[
    u \equiv x\;t_1 \cdots t_k
    \quad\text{(resp.\ }\quad
    t \equiv \lam{x_1 \cdots x_n}x\;u_1 \cdots u_k).
  \]
  That is, $\beta$-redex cannot exist in $t$ if $t$ is normal.
\end{frame}

\begin{frame}{Soundness and completeness of the inductive characterisation}
  A term $t$ has no $\beta$-reduction if and only if $t$ is normal:
  \begin{lemma}
    \begin{description}
      \item[Soundness] If $\normal\;t$ (resp.\ $\neutral\;t$) is derivable, then $t$ is in normal form.
      \item[Completeness]
        If $t$ is in normal form, then $\normal\;t$ is derivable.
    \end{description}
  \end{lemma}
  \begin{proof}[Proof sketch]
    \begin{description}
      \item[Soundness] By mutual induction on the derivation of
    $\normal\;t$ and $\neutral\;t$.
      \item[Completeness] By induction on the formation of $t$.
    \end{description}
  \end{proof}

  
\end{frame}

\begin{frame}{Progress}
  \begin{theorem}
    If $\Gamma \vdash t : A$ is derivable, then $\normal\;t$ or there is $u$ with $t \onereduce u$.
  \end{theorem}
  \begin{proof}[Proof sketch]
    By induction on the derivation of $\Gamma \vdash t : A$. 
  \end{proof}
\end{frame}


\begin{frame}{Weak normalisation}
  \begin{definition}
    $t$ is \emph{weakly normalising} denoted by $t\downarrow$ if 
    \begin{multicols}{2}
      \begin{prooftree}
        \AXC{$\normal\; t$}
        \UIC{$t \downarrow$}
      \end{prooftree}
      \columnbreak
      \begin{prooftree}
        \AXC{$t \onereduce u$}
        \AXC{$u \downarrow$}
        \BIC{$t \downarrow$}
      \end{prooftree}
    \end{multicols}
  \end{definition}
  That is, $t$ is weakly normalising if there is a sequence
  \[
    t \onereduce t_1 \onereduce t_2 \onereduce \dots u \notonereduce 
  \]
  \begin{theorem}[Weak normalisation]
    Every term $t$ with $\Gamma \vdash t : A$ is weakly normalising.
  \end{theorem}
\end{frame}

\begin{frame}{Strong normalisation}
  \begin{definition}
    $t$ is \emph{strongly normalising} denoted by $t \Downarrow$ if 
    \begin{prooftree}
      \AXC{$\forall u.\, (t \onereduce u \implies u \Downarrow)$}
      \UIC{$u \Downarrow$}
    \end{prooftree}
    
  \end{definition}
  Intuitively, \emph{strong normalisation} says every sequence
    \[
      t \onereduce t_1 \onereduce t_2 \cdots
    \]
  terminates, but the definition builds the sequence backwards. 

  \begin{theorem}
    Every term $t$ with $\Gamma \vdash t : A$ is strongly normalising.
  \end{theorem}
\end{frame}

%\begin{frame}{Typed Conversion}
%
%
%\end{frame}

%\begin{frame}{Definability}
%%Within simply typed lambda calculus, we wonder how many computable functions can
%%be defined, as it excludes non-terminating $\lambda$-terms. 
%  A function $f\colon \mathbb{N}^k \to \mathbb{N}$ is called
%  \alert{\emph{$\lambda_\to$-definable}} if there is a $\lambda$-term $F$ of
%  type $\nat \to \nat \to \dots \nat \to \nat$ such that
%  \[
%    F\,\bc_{n_1}\ldots\bc_{n_k} \reduce \bc_{f(n_1, \dots, n_k)}
%  \]
%  for every sequence $(n_1, n_2, \ldots, n_k) \in \mathbb{N}^k$.
%  Diagrammatically, 
%
%\[
%  \xymatrix{
%    (n_1, n_2, \ldots, n_k) \ar@{|->}[rr] \ar@{|->}[d]_{(\bc_{-})^k} & & f(n_1,
%    n_2, \ldots, n_k) \ar@{|->}[d]^{\bc_{-}}\\
%    (\bc_{n_1}, \bc_{n_2}, \ldots, \bc_{n_k}) \ar@{|->}[rr] & & 
%    F\;\bc_{n_1}\; \bc_{n_2}\; \ldots\;\bc_{n_k}
%    = \bc_{f(n_1, n_2, \ldots, n_k)}
%  }
%\]
%\end{frame}
%\begin{frame}{The limit of $\lambda_\to$}
%\begin{theorem}
%  The $\lambda_\to$-definable functions are the class of functions
%  of the form $f\colon \mathbb{N}^k \to \mathbb{N}$ closed under 
%  compositions
%  which contains
%  \begin{itemize}
%    \item the constant functions,
%    \item projections,
%    \item additions,
%    \item multiplications,
%    \item and the conditional 
%  \[
%    \mathrm{ifz}(n_0, n_1, n_2) = 
%    \begin{cases}
%      n_1 & \text{if } n_0 = 0\\
%      n_2 & \text{otherwise.}
%    \end{cases}
%  \]
%  \end{itemize}
%\end{theorem}
%\end{frame}
\section{Extensions to Simply Typed \texorpdfstring{$\lambda$}{λ}-Calculus}


\begin{frame}{General Recursion: Static}
  Self-applicative term cannot be typed in simply typed $\lambda$-calculus.
  E.g., 
  \[
    \lambda x.\, x\,x
  \]
  cannot be typed, since $A \to A$ is not equal to $A$.
  Hence, the $Y$-combinator in untyped $\lambda$-calculus cannot be typed.

  A construct is introduced explicitly for general recursion:

  Let $\Lambda_{\fix}(V)$ be the set of terms defined with an additional construct:
  \begin{description}
    \item[fixpoint] $\fix\,f.\, t$ is a term in $\Lambda_{\fix}(V)$, if $t \in \Lambda_{\fix}(V)$ and $f \in V$
  \end{description}

  An additional typing rule is added to simply typed $\lambda$-calculus:
  \begin{prooftree}
    \AXC{$\Gamma, f : A \vdash_{i} t : A$}
    \UIC{$\Gamma \vdash_{i} \fix\,f.\,t : A$}
  \end{prooftree}
\end{frame}

\begin{frame}{General Recursion: Dynamic}
  $\beta$-reduction for the general recursion $\fix$ is extended with the relation
  \begin{align*}
    \fix\,x.\,t & \onereduce t\subst{\fix\,x.\,t}{x}
  \end{align*}
  \vfill
  A term which never terminates can be defined easily.
  \begin{align*}
             & \fix\,x.\,x & \onereduce x\subst{\fix\,x.\,x}{x} \\
    \equiv{} & \fix\,x.\,x & \onereduce x\subst{\fix\,x.\,x}{x} \\
    \equiv{} & \fix\,x.\,x & \onereduce x\subst{\fix\,x.\,x}{x} \\
    \equiv{} & \dots
  \end{align*}
  \vfill
  Other notions such as $=_\alpha, \reduce$, and $\FV$ are extended similarly.
\end{frame}

%\subsection{System T}
\begin{frame}{Natural numbers: Static}
  While Church numerals can have multiple types $\nat_A$, for any $A$, we extend the calculus with a single type of natural numbers instead.

  Let $\Lambda_{\fix, \N}(V)$ be the set of terms defined with additional constructs:
  \begin{itemize}
    \item $\zero$ is a term in $\Lambda_{\fix, \N}(V)$
    \item $\suc(t)$ is a term in $\Lambda_{\fix, \N}(V)$ if $t$ is
    \item $\ifz(t; x.\,u)$ is a term in $\Lambda_{\fix, \N}(V)$ if $t, u \in \Lambda_{\fix, \N}(V)$ and $x \in V$
  \end{itemize}
  with additional typing rules
  \begin{multicols}{2}
    \begin{prooftree}
      \AXC{$\vphantom{\Gamma}$}
      \UIC{$\Gamma \vdash \zero : \mathbb{N}$}
    \end{prooftree}
    \begin{prooftree}
      \AXC{$\Gamma \vdash t : \mathbb{N}$}
      \UIC{$\Gamma \vdash \suc(t) : \mathbb{N}$}
    \end{prooftree}
  \end{multicols}

  \begin{prooftree}
    \AXC{$\Gamma \vdash v : \mathbb{N}$}
    \AXC{$\Gamma \vdash t : A$}
    \AXC{$\Gamma, x : \mathbb{N} \vdash u : A$}
    \TrinaryInfC{$\Gamma \vdash \ifz(t; x.\,u)\;v : A$}
  \end{prooftree}
  The third rule is akin to pattern matching on natural numbers.
\end{frame}

\begin{frame}{Natural numbers: Dynamic}
  $\beta$-reduction for natural numbers is extended with two rules:
  \begin{align*}
    \ifz(t; x.\, u)\;\zero   & \onereduce t \\
    \ifz(t; x.\, u)\;\suc(n) & \onereduce u\subst{n}{x}
  \end{align*}
\end{frame}

\begin{frame}{Natural numbers: Exercise}
  Define the predecessor of natural numbers as a program
  \[
    \pred : \mathbb{N} \to \mathbb{N}.
  \]

  Evaluate the following terms to their normal forms.
  \begin{enumerate}
    \item $\pred\;\zero$
    \item $\pred\;(\suc\;(\suc\;(\suc\;\zero)))$
  \end{enumerate}
\end{frame}

\begin{frame}{Boolean values: Exercise}
  Extend simply typed $\lambda$-calculus $\Lambda_{\fix, \N}(V)$ further with a type of Boolean values. 
  \begin{enumerate}
    \item What term constructs are needed?
    \item What typing rules should be added?
    \item How $\beta$-reduction should be updated?
    \item Define Boolean operations, i.e. conjunction, disjunction, and negation, in this extension.
  \end{enumerate}

\end{frame}
%include numerals in the generation of syntax:
%  \begin{multicols}{2}
%    \begin{prooftree}
%      \AXC{$x \in V$}
%      \UIC{$x \in \Lambda_{\T, \nat}$}
%    \end{prooftree}
%    \begin{prooftree}
%      \AXC{\phantom{$n \in \mathbb{N}$}}
%      \UIC{$\zero\in \Lambda_{\T, \nat}$}
%    \end{prooftree}
%    \begin{prooftree}
%      \AXC{$M \in \Lambda_{\T, \nat}$}
%      \UIC{$\suc\;M\in \Lambda_{\T, \nat}$}
%    \end{prooftree}
%    \begin{prooftree}
%      \AXC{$M \in \Lambda_{\T,\nat}$}
%      \AXC{$N \in \Lambda_{\T,\nat}$}
%      \BIC{$(M\, N) \in \Lambda_{\T, \nat}$}
%    \end{prooftree}
%    \begin{prooftree}
%      \AXC{$M \in \Lambda_{\T, \nat}$}
%      \AXC{$x \in V$}
%      \AXC{$\tau \in \mathbb{T}$}
%      \TIC{$\lambda (x:\tau).\; M \in \Lambda_{\T, \nat}$}
%    \end{prooftree}
%  \end{multicols}
%\noindent and let $\mathbb{G} \defeq \{\nat\}$ with an addition set of typing
%rules:
%\begin{multicols}{2}
%\begin{prooftree}
%  \AXC{\phantom{$\Gamma$}}
%  \UIC{$\Gamma \vdash \zero : \nat$}
%\end{prooftree}
%\begin{prooftree}
%  \AXC{$\Gamma \vdash M : \nat$}
%  \UIC{$\Gamma \vdash \suc\:M : \nat$}
%\end{prooftree}
%\end{multicols}
%\noindent In this way, we can derive that every $\underline{n} \defeq
%\suc^n\;\zero$ has the type of $\nat$:
%\[
%  \Gamma \vdash \underline{n} : \nat
%\]
%
%However, it still remains to add primitive operations such as addition,
%multiplication, division, and so on with a proper set of reduction rules. As we
%have seen that recursions on natural numbers in untyped lambda calculus can be
%defined with the fixpoint operator. Something similar can be done:
%\begin{definition}
%  In addition to typing rules introduced so far, we add the following
%  \begin{prooftree}
%    \AXC{$\Gamma \vdash M : \nat$}
%    \AXC{$\Gamma \vdash M_0 : \sigma$}
%    \AXC{$\Gamma \vdash F : \nat \to (\sigma \to \sigma)$}
%    \TIC{$\Gamma \vdash \fold\;M_0\;F : \nat \to \sigma $}
%  \end{prooftree}
%  with additional reductions
%  \[
%    \fold\;M_0\;F\;\zero
%    \xrightarrow{\beta}
%    M_0
%    \quad\text{and}\quad
%    \fold\;M_0\;F\;(\suc\;M)
%    \xrightarrow{\beta}
%    F\;M\;(\fold\;F\;M).
%  \]
%\end{definition}
%\begin{definition}
%  \todo[inline]{define addition, multiplication using $\fold$}
%\end{definition}
%\begin{frame}{Remark}
%With the decidability of type checking, Preservation Theorem, Progress Theorem
%(a well-typed term is either a ``value'' or a reducible term), and the strong
%normalisation, we actually have exhibited a decidable evaluator of simply typed
%lambda calculus that always reduce a well-typed term of type~$\sigma$ to a value
%of type~$\sigma$.
%\end{frame}

\begin{frame}{Homework}
  \begin{enumerate}
    \item (2.5\%) Complete the proof of the Preservation Theorem.
    \item (5\%) Show the Progress Theorem.
    \item (2.5\%) Show that if $t$ is in normal form then $\normal\;t$ is derivable.
    \item (5\%) Extend $\Lambda_{\fix, \N}(V)$ further with product types $A \times B$, for any $A$ and $B$ where additional constructs should include pairs $(t, u)$ and a construct to pattern match on a pair.

  \end{enumerate}
  
\end{frame}

%\begin{frame}{References}
%
%\bibliographystyle{amsalpha}
%\bibliography{library} 
%
%\end{frame}

\appendix

\end{document}
